Sustainable Development Goal~11.7 calls for universal access to green and public
space, yet accessibility is almost always quantified as a static snapshot that
says nothing about how access degrades when the street network is disrupted. We
introduce a \emph{topological resilience metric} for accessibility: the
Wasserstein distance between the persistent-homology diagram of a city's
green-space accessibility field before and after network disruption, traced as a
function of disruption intensity~$\rho$. The area under this degradation curve
and its critical transition $\rho^{*}$ summarise how robustly a place keeps its
population close to nature. Applying the metric to five morphologically
contrasting cities (İstanbul, Barcelona, Amsterdam, Bogotá, Phoenix) at the
level of $2{,}603$ intra-urban districts, we find that street-network morphology
is a strong, statistically significant predictor of green-space access
resilience: in a district-level regression with city fixed effects, street
circuity, orientation entropy, mean segment length and intersection density are
all significant ($p<0.01$; three at $p<10^{-8}$). The relationship is robust: the
three strongest effects keep their sign and significance under both random and
targeted disruption and across district resolutions spanning a $28\times$ range of
unit counts. The signal lives in dimension-0 homology---the merging of well-served
regions---rather than in enclosed ``access deserts,'' a finding that only emerges
on real street networks.
